\part{统计学基础}

\chapter{概率论回顾}

\section{随机变量与分布}

\subsection{离散随机变量}

\begin{definition}[概率质量函数]
对于离散随机变量 $X$，其概率质量函数 (PMF) $p(x)$ 满足：
\begin{equation}
    p(x) = P(X = x),\quad \sum_{x} p(x) = 1.
\end{equation}
\end{definition}

\begin{definition}[期望]
离散随机变量 $X$ 的期望定义为：
\begin{equation}
    \E[X] = \sum_{x} x\,p(x).
\end{equation}
\end{definition}

\begin{definition}[方差]
随机变量 $X$ 的方差定义为：
\begin{equation}
    \Var(X) = \E\big[(X - \E[X])^2\big] = \E[X^2] - (\E[X])^2.
\end{equation}
\end{definition}

\subsection{连续随机变量}

\begin{definition}[概率密度函数]
对于连续随机变量 $X$，其概率密度函数 (PDF) $f(x)$ 满足：
\begin{equation}
    P(a \le X \le b) = \int_{a}^{b} f(x)\,dx,\quad \int_{-\infty}^{\infty} f(x)\,dx = 1.
\end{equation}
\end{definition}

\begin{definition}[期望与方差]
连续随机变量 $X$ 的期望与方差：
\begin{align}
    \E[X]    &= \int_{-\infty}^{\infty} x f(x)\,dx, \\
    \Var(X)  &= \int_{-\infty}^{\infty} (x - \E[X])^2 f(x)\,dx.
\end{align}
\end{definition}

\subsection{常见分布}

\begin{enumerate}[label=(\arabic*)]
    \item \textbf{二项分布} $X \sim \mathrm{Binomial}(n, p)$：
    \begin{equation}
        P(X = k) = \binom{n}{k} p^k (1-p)^{n-k},\quad \E[X] = np,\; \Var(X) = np(1-p).
    \end{equation}
    \item \textbf{泊松分布} $X \sim \mathrm{Poisson}(\lambda)$：
    \begin{equation}
        P(X = k) = \frac{\lambda^k e^{-\lambda}}{k!},\quad \E[X] = \lambda,\; \Var(X) = \lambda.
    \end{equation}
    \item \textbf{正态分布} $X \sim \mathcal{N}(\mu, \sigma^2)$：
    \begin{equation}
        f(x) = \frac{1}{\sqrt{2\pi\sigma^2}} \exp\!\left(-\frac{(x-\mu)^2}{2\sigma^2}\right).
    \end{equation}
    \item \textbf{多元正态分布} $\mathbf{X} \sim \mathcal{N}(\bm{\mu}, \bm{\Sigma})$：
    \begin{equation}
        f(\mathbf{x}) = \frac{1}{(2\pi)^{d/2}|\bm{\Sigma}|^{1/2}}
        \exp\!\left(-\frac{1}{2}(\mathbf{x}-\bm{\mu})^\top \bm{\Sigma}^{-1}(\mathbf{x}-\bm{\mu})\right).
    \end{equation}
\end{enumerate}

\section{条件概率与独立性}

\begin{definition}[条件概率]
事件 $A$ 在事件 $B$ 发生条件下的概率：
\begin{equation}
    P(A \mid B) = \frac{P(A \cap B)}{P(B)},\quad P(B) > 0.
\end{equation}
\end{definition}

\begin{theorem}[贝叶斯公式]
\begin{equation}
    P(B_k \mid A) = \frac{P(A \mid B_k) P(B_k)}{\sum_{i} P(A \mid B_i) P(B_i)}.
\end{equation}
\end{theorem}

\begin{definition}[独立性]
若 $P(A \cap B) = P(A)P(B)$，则事件 $A$ 与 $B$ 独立。
\end{definition}

\section{矩与特征}

\begin{definition}[协方差]
随机变量 $X$ 与 $Y$ 的协方差：
\begin{equation}
    \Cov(X, Y) = \E\big[(X - \E[X])(Y - \E[Y])\big].
\end{equation}
\end{definition}

\begin{definition}[相关系数]
\begin{equation}
    \rho_{XY} = \frac{\Cov(X, Y)}{\sqrt{\Var(X)\,\Var(Y)}} \in [-1, 1].
\end{equation}
\end{definition}

\subsection{Pearson 与 Spearman 相关系数}

上面定义的是理论上的（总体）相关系数 $\rho$。实际中，我们需要用样本数据估计它，
常用的有两种：Pearson 和 Spearman。

\begin{definition}[Pearson 相关系数 $r$]
样本 Pearson 相关系数度量的是\textbf{线性}相关程度：
\begin{equation}
    r = \frac{\sum_{i=1}^{n}(x_i - \bar{x})(y_i - \bar{y})}
             {\sqrt{\sum_{i=1}^{n}(x_i - \bar{x})^2}\;
              \sqrt{\sum_{i=1}^{n}(y_i - \bar{y})^2}}.
\end{equation}
\end{definition}

\begin{definition}[Spearman 秩相关系数 $\rho_s$]
将原始数据 $x_i, y_i$ 分别替换为它们在各自序列中的\textbf{排名}（秩），
然后对排名计算 Pearson 相关系数，即得 Spearman 相关系数。

等价公式（无结时）：
\begin{equation}
    \rho_s = 1 - \frac{6\sum d_i^2}{n(n^2-1)},
\end{equation}
其中 $d_i$ 为第 $i$ 对数据的 $x$ 排名与 $y$ 排名之差。
\end{definition}

\begin{definition}[Kendall 秩相关系数 $\tau$]
Kendall's $\tau$ 同样基于排名，但计算逻辑不同——它统计的是\textbf{一致对}与\textbf{不一致对}的数量。

对任意两对观测 $(x_i, y_i)$ 和 $(x_j, y_j)\;(i < j)$：
\begin{itemize}
    \item 若 $(x_i - x_j)(y_i - y_j) > 0$：\textbf{一致对}（两变量同方向变化）
    \item 若 $(x_i - x_j)(y_i - y_j) < 0$：\textbf{不一致对}（两变量反方向变化）
\end{itemize}

令 $n_c$ 为一致对数，$n_d$ 为不一致对数，则：
\begin{equation}
    \tau = \frac{n_c - n_d}{\frac{1}{2}n(n-1)}.
\end{equation}

$\tau = 1$：所有对都一致（完美正相关）；$\tau = -1$：所有对都不一致（完美负相关）。
\end{definition}

\medskip
\textbf{三者对比：}

\begin{center}
\renewcommand{\arraystretch}{1.4}
\small
\begin{tabular}{p{2.2cm}|p{3.5cm}|p{3.5cm}|p{3.5cm}}
\toprule
& \textbf{Pearson $r$} & \textbf{Spearman $\rho_s$} & \textbf{Kendall $\tau$} \\
\midrule
度量什么   & 线性相关程度 & 单调相关程度 & 排序一致程度 \\
计算方式   & 原始值 & 原始值的\textbf{排名} & 一致对 vs 不一致对 \\
对异常值   & \textbf{敏感} & \textbf{稳健} & \textbf{最稳健} \\
对结（ties）& 无影响 & 需要修正公式 & 天然处理结 \\
样本效率   & 最高 & 中等 & 较低（需更多样本） \\
适用场景   & 近似正态、线性 & 非正态、单调 & 小样本、多结、定序数据 \\
\midrule
\end{tabular}
\end{center}

\textbf{使用规则：}
\begin{itemize}
    \item 先画散点图。如果散点大致沿一条直线分布 → 用 Pearson；
    \item 如果散点呈曲线但有单调趋势 → 用 Spearman；
    \item 数据有严重离群值或明显不正态 → Spearman 或 Kendall；
    \item 样本量小（$n<30$）、有很多重复值（结） → 优先 Kendall；
    \item \textbf{最佳实践：} 三者都算。若结论一致（如都显著且方向相同），说明相关性结论\textbf{稳健}，可放心报告 Pearson。
\end{itemize}

\medskip
\textbf{$r$ 与 $r^2$ 的区别：}

\begin{center}
\renewcommand{\arraystretch}{1.3}
\begin{tabular}{c|c|c}
\toprule
& \textbf{$r$（相关系数）} & \textbf{$r^2$（决定系数）} \\
\midrule
范围   & $[-1, 1]$（有方向） & $[0, 1]$（无方向） \\
含义   & 方向和强度 & 变异解释比例 \\
例子   & $r = -0.5$ 表示中等负相关 & $r^2 = 0.25$ 表示 $X$ 解释了 $Y$ 25\% 的变异 \\
关系   & \multicolumn{2}{c}{$r^2 = r \times r$（仅在一元回归中等价于 $R^2$）} \\
\bottomrule
\end{tabular}
\end{center}

注意：
\begin{itemize}
    \item $r^2$ 只在\textbf{一元线性回归}中等价于第 9 章定义的 $R^2$；
    \item 多元回归的 $R^2$ 是所有自变量共同解释的比例，不是某个 $r$ 的平方。
\end{itemize}

